% ---
% titre: Calcul mental
% theme: NO
% chapitre: Nombres naturels et décimaux
% sous_chapitre: PPMC et PGDC
% niveau: [9e]
% type: EVAL
% tags: [c-nombres-premiers, c-decomposition-facteurs-premiers, c-multiple, c-diviseur, c-ppmc,  p-question-TS]
% difficulte: 1
% corrige: true
% ---

\question
Détermine:

\begin{parts} 
\vspace{10pt}\part[3]
le ppmc de 16 et 40

\vspace{10pt}{\footnotesize\raggedleft Espace pour ta démarche et tes calculs\par}
	\begin{solutionorgrid}[110pt]
		\begin{minipage}[t]{0.1\linewidth}
		\vspace{0pt}
		\begin{tabularx}{\linewidth}{>{\raggedleft\arraybackslash}X|X}
		16&2\\
		8&2\\
		4&2\\
		2&2\\
		1&\\
		\end{tabularx}
		\end{minipage}%
		\hfill
		\begin{minipage}[t]{0.1\linewidth}
		\vspace{0pt}
		\begin{tabularx}{\linewidth}{>{\raggedleft\arraybackslash}X|X}
		40&2\\
		20&2\\
		10&2\\
		5&5\\
		1&\\
		\end{tabularx}
		\end{minipage}%
		\hfill
		\begin{minipage}[t]{0.6\linewidth}
		\vspace{0pt}
		$16 = 2^4$ \quad / \quad $40 = 2^3 \cdot 5$
		
		ppmc(16;40) $= 2^4 \cdot 5 = 80$
		
		\vspace{20pt}
		1pt pour chaque décomposition ou liste de multiples
		
		1pt pour le ppmc
		\end{minipage}%
	\end{solutionorgrid}

\vspace{10pt}\part[3]
le ppmc de 39 et 21

\vspace{10pt}{\footnotesize\raggedleft Espace pour ta démarche et tes calculs\par}
	\begin{solutionorgrid}[110pt]
		\begin{minipage}[t]{0.1\linewidth}
		\vspace{0pt}
		\begin{tabularx}{\linewidth}{>{\raggedleft\arraybackslash}X|X}
		39&3\\
		13&13\\
		1&\\
		\end{tabularx}
		\end{minipage}%
		\hfill
		\begin{minipage}[t]{0.1\linewidth}
		\vspace{0pt}
		\begin{tabularx}{\linewidth}{>{\raggedleft\arraybackslash}X|X}
		21&3\\
		7&7\\
		1&\\
		\end{tabularx}
		\end{minipage}%
		\hfill
		\begin{minipage}[t]{0.6\linewidth}
		\vspace{0pt}
		$39 = 3 \cdot 13$ \quad / \quad $21 = 3 \cdot 7$
		
		ppmc(39;21) $= 3 \cdot 7 \cdot 13 = 273$
		
		\vspace{20pt}
		1pt pour chaque décomposition ou liste de multiples
		
		1pt pour le ppmc
		\end{minipage}%
	\end{solutionorgrid}

\vspace{10pt}\part[3]
le pgdc de 12 et 14

\vspace{10pt}{\footnotesize\raggedleft Espace pour ta démarche et tes calculs\par}
	\begin{solutionorgrid}[110pt]
		\begin{minipage}[t]{0.1\linewidth}
		\vspace{0pt}
		\begin{tabularx}{\linewidth}{>{\raggedleft\arraybackslash}X|X}
		12&2\\
		6&2\\
		3&3\\
		1&\\
		\end{tabularx}
		\end{minipage}%
		\hfill
		\begin{minipage}[t]{0.1\linewidth}
		\vspace{0pt}
		\begin{tabularx}{\linewidth}{>{\raggedleft\arraybackslash}X|X}
		14&2\\
		7&7\\
		1&\\
		\end{tabularx}
		\end{minipage}%
		\hfill
		\begin{minipage}[t]{0.6\linewidth}
		\vspace{0pt}
		$12 = 2^2 \cdot 3$ \quad / \quad $14 = 2 \cdot 7$
		
		pgdc(12;14) $= 2$
		
		\vspace{20pt}
		1pt pour chaque décomposition ou liste de diviseurs
		
		1pt pour le ppmc
		\end{minipage}%
	\end{solutionorgrid}

\vspace{10pt}\part[3]
le pgdc de 120 et 225

\vspace{10pt}{\footnotesize\raggedleft Espace pour ta démarche et tes calculs\par}
	\begin{solutionorgrid}[110pt]
		\begin{minipage}[t]{0.1\linewidth}
		\vspace{0pt}
		\begin{tabularx}{\linewidth}{>{\raggedleft\arraybackslash}X|X}
		120&2\\
		60&2\\
		30&2\\
		15&3\\
		5&5\\
		1&\\
		\end{tabularx}
		\end{minipage}%
		\hfill
		\begin{minipage}[t]{0.1\linewidth}
		\vspace{0pt}
		\begin{tabularx}{\linewidth}{>{\raggedleft\arraybackslash}X|X}
		225&3\\
		75&3\\
		25&5\\
		5&5\\
		1&\\
		\end{tabularx}
		\end{minipage}%
		\hfill
		\begin{minipage}[t]{0.6\linewidth}
		\vspace{0pt}
		$120 = 2^3 \cdot 3 \cdot 5$ \quad / \quad $225 = 3^2 \cdot 5^2$
		
		pgdc(120;225) $= 3 \cdot 5 = 15$
		
		\vspace{20pt}
		1pt pour chaque décomposition ou liste de diviseurs
		
		1pt pour le ppmc
		\end{minipage}%
	\end{solutionorgrid}
\end{parts}

